% ===========================================================================
%  Chapitre 4 — Primitives et calcul intégral
%  PE 2026, composante ANALYSE
% ===========================================================================
\bmtchapitre{Primitives et calcul intégral}
  {Calculer une primitive et une intégrale ; intégrer par parties et par
   changement de variable affine ; calculer une aire, une valeur moyenne et le
   volume d'un solide de révolution.}
  {Analyse \textbullet\ environ 12 heures}

Jusqu'ici, une fonction étant donnée, vous calculiez sa dérivée. Ce chapitre
prend le problème à l'envers : une dérivée étant donnée, retrouver la fonction
dont elle provient. L'opération n'a rien d'automatique — dériver est une
mécanique, primitiver est une reconnaissance.

Et il se produit alors quelque chose d'inattendu. Cette question purement
algébrique, remonter d'une dérivée à sa fonction, donne la réponse à une
question de géométrie que personne n'aurait crue liée : quelle est l'aire
sous une courbe. C'est le résultat le plus surprenant du programme de
terminale, et le plus fécond : il permettra de calculer des aires, des
volumes, des valeurs moyennes, et de majorer des suites d'intégrales dont on
ne sait rien par ailleurs.

% ---------------------------------------------------------------------------
\begin{revision}
Uniquement des dérivées — c'est le sens que ce chapitre va renverser.

\begin{enumerate}[label=\textbf{\arabic*.}, leftmargin=*, itemsep=0.35em]
  \item Dériver : $x \mapsto x^{5}$, $x \mapsto \sqrt{x}$,
        $x \mapsto \dfrac{1}{x^{2}}$, $x \mapsto \sin(2x)$.
  \item Dériver $x \mapsto (3x+1)^{4}$, puis $x \mapsto \sqrt{2x-5}$.
  \item Rappeler la formule de dérivation d'un produit.
  \item Rappeler la formule $\left(u^{n}\right)'$ pour $n$ entier.
  \item Quelle est l'aire d'un rectangle de largeur $b-a$ et de hauteur $h$ ?
  \item Quel est le volume d'un cylindre de rayon $r$ et de hauteur $h$ ?
\end{enumerate}
\end{revision}

\begin{automatismes}
Sans calculatrice, sans brouillon.

\begin{enumerate}[label=\textbf{\alph*)}, leftmargin=*, itemsep=0.25em]
  \item Dérivée de $x \mapsto \dfrac{x^{4}}{4}$
  \item Dérivée de $x \mapsto 2\sqrt{x}$
  \item Dérivée de $x \mapsto -\dfrac{1}{x}$
  \item Dérivée de $x \mapsto -\dfrac{\cos(3x)}{3}$
  \item Dérivée de $x \mapsto \dfrac{(2x+1)^{5}}{10}$
  \item Une fonction dont la dérivée est $x \mapsto x$
  \item Une fonction dont la dérivée est $x \mapsto \cos x$
  \item Une fonction dont la dérivée est $x \mapsto 3x^{2}$
\end{enumerate}
\end{automatismes}

\begin{corrige}[automatismes]
a) $x^{3}$ \quad b) $\dfrac{1}{\sqrt{x}}$ \quad c) $\dfrac{1}{x^{2}}$ \quad
d) $\sin(3x)$ \quad e) $(2x+1)^{4}$ \quad f) $\dfrac{x^{2}}{2}$ \quad
g) $\sin x$ \quad h) $x^{3}$. Les quatre dernières annoncent exactement ce que
le chapitre va formaliser.
\end{corrige}

% ===========================================================================
\section{Les primitives}

\begin{definition}[primitive]
Soit $f$ une fonction définie sur un intervalle $I$. On appelle
\textbf{primitive} de $f$ sur $I$ toute fonction $F$ dérivable sur $I$ telle
que
\[
  F'(x) = f(x) \qquad \text{pour tout } x \in I .
\]
\end{definition}

\begin{theoreme}[existence]
Toute fonction continue sur un intervalle $I$ admet des primitives sur $I$.
\end{theoreme}

\begin{theoreme}[deux primitives diffèrent d'une constante]
Si $F$ est une primitive de $f$ sur $I$, alors les primitives de $f$ sur $I$
sont exactement les fonctions $x \mapsto F(x)+k$, où $k$ décrit $\R$.
\end{theoreme}

\begin{demonstration}
Si $G$ est une autre primitive, la fonction $H = G-F$ vérifie
$H' = G'-F' = f-f = 0$ sur l'intervalle $I$. Une fonction de dérivée nulle sur
un intervalle y est constante — c'est une conséquence directe du théorème des
accroissements finis, appliqué entre deux points quelconques $a$ et $b$ de
$I$ : $H(b)-H(a) = (b-a)H'(c) = 0$. Donc $G = F+k$. Réciproquement, toute
fonction $F+k$ a bien $f$ pour dérivée.
\end{demonstration}

\begin{avertissement}
L'énoncé ci-dessus vaut sur un \emph{intervalle}. Sur une réunion de deux
intervalles disjoints, les constantes peuvent différer d'un morceau à l'autre.
Écrire « la primitive de $\frac{1}{x^{2}}$ sur $\Rs$ est $-\frac{1}{x}+k$ »
est donc imprécis : il faut traiter séparément $\left]-\infty;0\right[$ et
$\left]0;+\infty\right[$.
\end{avertissement}

\subsection{Le tableau des primitives usuelles}

\begin{center}
\renewcommand{\arraystretch}{1.6}
\begin{tabular}{@{}p{4.3cm} p{4.3cm} p{3.6cm}@{}}
\toprule
\textsf{\footnotesize\bfseries FONCTION $f$} &
\textsf{\footnotesize\bfseries UNE PRIMITIVE $F$} &
\textsf{\footnotesize\bfseries SUR} \\
\midrule
$x \mapsto a$ (constante)        & $x \mapsto ax$                       & $\R$ \\
$x \mapsto x^{n}$, $n \in \N$    & $x \mapsto \dfrac{x^{n+1}}{n+1}$     & $\R$ \\
$x \mapsto \dfrac{1}{x^{2}}$     & $x \mapsto -\dfrac{1}{x}$            & tout intervalle de $\Rs$ \\
$x \mapsto \dfrac{1}{\sqrt{x}}$  & $x \mapsto 2\sqrt{x}$                & $\Rps$ \\
$x \mapsto \cos x$               & $x \mapsto \sin x$                   & $\R$ \\
$x \mapsto \sin x$               & $x \mapsto -\cos x$                  & $\R$ \\
\midrule
$u'u^{n}$, $n \neq -1$           & $\dfrac{u^{n+1}}{n+1}$               & où $u$ est dérivable \\
$\dfrac{u'}{\sqrt{u}}$           & $2\sqrt{u}$                          & où $u > 0$ \\
$u'\cos u$                       & $\sin u$                             & où $u$ est dérivable \\
$u' \times (v'\circ u)$          & $v \circ u$                          & \\
\bottomrule
\end{tabular}
\end{center}

\bmtmarge{La seconde moitié du tableau est la plus utile : elle se lit comme
les formules de dérivation prises à l'envers. Reconnaître un $u'u^{n}$ dans
une expression compliquée est la compétence centrale du chapitre.}

\begin{exemple}[reconnaître la forme $u'u^{n}$]
Cherchons une primitive de $f(x) = 3x^{2}\left(x^{3}+1\right)^{4}$ sur $\R$.

On pose $u(x) = x^{3}+1$, donc $u'(x) = 3x^{2}$. L'expression est exactement
$u'u^{4}$, dont une primitive est $\dfrac{u^{5}}{5}$. Donc
\[
  F(x) = \frac{\left(x^{3}+1\right)^{5}}{5}
\]
convient, et l'on vérifie en dérivant.

\smallskip
Si l'énoncé avait donné $f(x) = x^{2}\left(x^{3}+1\right)^{4}$, il aurait fallu
faire apparaître le facteur manquant :
$f(x) = \dfrac{1}{3}\times 3x^{2}\left(x^{3}+1\right)^{4}$, et donc
$F(x) = \dfrac{\left(x^{3}+1\right)^{5}}{15}$.
\end{exemple}

\begin{propriete}[primitive prenant une valeur donnée en un point]
Soit $f$ continue sur $I$, $x_{0} \in I$ et $y_{0}$ un réel. Il existe une
\emph{unique} primitive $F$ de $f$ sur $I$ telle que $F(x_{0}) = y_{0}$.
\end{propriete}

\begin{exemple}[fixer la constante]
Déterminons la primitive $F$ de $f(x) = 2x-3$ sur $\R$ telle que $F(1) = 4$.

Les primitives de $f$ sont les $x \mapsto x^{2}-3x+k$. La condition impose
$1-3+k = 4$, soit $k = 6$. Donc $F(x) = x^{2}-3x+6$, et c'est la seule.
\end{exemple}

\begin{entrainement}[primitives]
\exo[\niveau{1}] Déterminer une primitive de $x \mapsto 4x^{3}-2x+5$ sur $\R$,
puis de $x \mapsto \dfrac{3}{x^{2}}$ sur $\Rps$.

\exo[\niveau{2}] Déterminer une primitive de
$x \mapsto \dfrac{2x}{\sqrt{x^{2}+1}}$ sur $\R$.

\exo[\niveau{2}] Déterminer une primitive de $x \mapsto \sin x \cos^{3} x$ sur
$\R$. \emph{On posera $u = \cos x$.}

\exo[\niveau{3}] Déterminer la primitive $F$ de
$x \mapsto \dfrac{x}{\left(x^{2}+1\right)^{2}}$ sur $\R$ telle que $F(0) = 1$.
\end{entrainement}

\begin{corrige}
\textbf{1.} $x^{4}-x^{2}+5x$ ; \ $-\dfrac{3}{x}$.

\textbf{2.} C'est $\dfrac{u'}{\sqrt{u}}$ avec $u = x^{2}+1$ : une primitive est
$2\sqrt{x^{2}+1}$.

\textbf{3.} Avec $u = \cos x$, $u' = -\sin x$, l'expression vaut $-u'u^{3}$,
dont une primitive est $-\dfrac{u^{4}}{4} = -\dfrac{\cos^{4}x}{4}$.

\textbf{4.} L'expression est $\dfrac{1}{2}u'u^{-2}$ avec $u = x^{2}+1$, de
primitive $-\dfrac{1}{2u} = \dfrac{-1}{2(x^{2}+1)}$. La condition $F(0)=1$
donne $-\tfrac12+k=1$, soit $k = \tfrac32$ :
$F(x) = \dfrac{-1}{2(x^{2}+1)} + \dfrac{3}{2}$.
\end{corrige}

% ===========================================================================
\section{L'intégrale}

\begin{definition}[intégrale d'une fonction continue]
Soit $f$ continue sur un intervalle $I$, et $a$, $b$ deux éléments de $I$. Si
$F$ est une primitive quelconque de $f$ sur $I$, on pose
\[
  \int_{a}^{b} f(x)\dx = F(b)-F(a),
\]
quantité que l'on note aussi $\bigl[F(x)\bigr]_{a}^{b}$.
\end{definition}

\begin{remarque}
Le résultat ne dépend pas de la primitive choisie : si l'on remplace $F$ par
$F+k$, la différence $\bigl(F(b)+k\bigr)-\bigl(F(a)+k\bigr)$ est inchangée. La
définition est donc légitime.
\end{remarque}

\begin{propriete}[règles de calcul]
Pour $f$ et $g$ continues sur $I$, $a$, $b$, $c$ dans $I$ et $\lambda$ réel :
\begin{align*}
  \int_{a}^{a} f(x)\dx &= 0, &
  \int_{b}^{a} f(x)\dx &= -\int_{a}^{b} f(x)\dx, \\[0.3em]
  \int_{a}^{b} \bigl[f(x)+g(x)\bigr]\dx &= \int_{a}^{b} f + \int_{a}^{b} g, &
  \int_{a}^{b} \lambda f(x)\dx &= \lambda \int_{a}^{b} f(x)\dx,
\end{align*}
et la \textbf{relation de Chasles} :
\[
  \int_{a}^{b} f(x)\dx = \int_{a}^{c} f(x)\dx + \int_{c}^{b} f(x)\dx .
\]
\end{propriete}

\begin{propriete}[positivité et ordre]
Si $a \leq b$ :
\begin{itemize}[leftmargin=1.4em, itemsep=0.18em]
  \item $f \geq 0$ sur $\intervalle{a}{b}$ entraîne $\displaystyle\int_{a}^{b} f(x)\dx \geq 0$ ;
  \item $f \leq g$ sur $\intervalle{a}{b}$ entraîne
        $\displaystyle\int_{a}^{b} f(x)\dx \leq \int_{a}^{b} g(x)\dx$ ;
  \item si $m \leq f(x) \leq M$ sur $\intervalle{a}{b}$, alors
        $m(b-a) \leq \displaystyle\int_{a}^{b} f(x)\dx \leq M(b-a)$.
\end{itemize}
\end{propriete}

\begin{methode}{encadrer une intégrale qu'on ne sait pas calculer}
La troisième propriété est celle qui tombe au baccalauréat, presque toujours
sur une suite $(I_{n})$ d'intégrales. Le mécanisme ne varie pas :

\begin{enumerate}[label=\textbf{\arabic*.}, leftmargin=*, itemsep=0.28em]
  \item encadrer la \emph{fonction} sous l'intégrale sur l'intervalle
        d'intégration, sans jamais chercher à la primitiver ;
  \item intégrer l'encadrement terme à terme, en respectant le sens des
        inégalités ;
  \item conclure sur la limite de $(I_{n})$ par le théorème des gendarmes.
\end{enumerate}

Notez qu'on ne calcule à aucun moment $I_{n}$ : c'est tout l'intérêt.
\end{methode}

% ===========================================================================
\section{Deux techniques de calcul}

\subsection{L'intégration par parties}

\begin{theoreme}[intégration par parties]
Soient $u$ et $v$ deux fonctions dérivables sur $\intervalle{a}{b}$, de
dérivées continues. Alors
\[
  \int_{a}^{b} u(x)v'(x)\dx
  = \bigl[u(x)v(x)\bigr]_{a}^{b} - \int_{a}^{b} u'(x)v(x)\dx .
\]
\end{theoreme}

\begin{demonstration}
La formule de dérivation d'un produit donne $(uv)' = u'v + uv'$. Les trois
fonctions étant continues sur $\intervalle{a}{b}$, on intègre cette égalité
entre $a$ et $b$ :
\[
  \bigl[u(x)v(x)\bigr]_{a}^{b} = \int_{a}^{b} u'(x)v(x)\dx + \int_{a}^{b} u(x)v'(x)\dx,
\]
puisque $uv$ est une primitive de $(uv)'$. Il ne reste qu'à isoler la seconde
intégrale.
\end{demonstration}

\begin{methode}{choisir $u$ et $v'$}
Tout le calcul dépend de ce choix, et le critère est simple : on prend pour
$u$ la fonction que la dérivation \emph{simplifie}, et pour $v'$ celle que
l'on sait primitiver sans peine.

Un polynôme multiplié par un cosinus ou un sinus : le polynôme est $u$, car sa
dérivée baisse d'un degré et finit par disparaître, tandis que le sinus et le
cosinus se primitivent indéfiniment. Un polynôme multiplié par une racine
carrée : le polynôme est encore $u$. Les chapitres 5 et 6 ajouteront deux
familles à cette liste, et le critère restera le même.
\end{methode}

\begin{exemple}[une intégration par parties menée en entier]
Calculons $\displaystyle\int_{0}^{\pi} x\sin x \dx$.

On pose $u(x) = x$ et $v'(x) = \sin x$, d'où $u'(x) = 1$ et
$v(x) = -\cos x$. La formule donne
\[
  \int_{0}^{\pi} x\sin x \dx
  = \bigl[-x\cos x\bigr]_{0}^{\pi} - \int_{0}^{\pi} (-\cos x)\dx
  = \bigl[-x\cos x\bigr]_{0}^{\pi} + \bigl[\sin x\bigr]_{0}^{\pi} .
\]
Or $\bigl[-x\cos x\bigr]_{0}^{\pi} = -\pi\cos\pi + 0 = \pi$ et
$\bigl[\sin x\bigr]_{0}^{\pi} = 0-0 = 0$. Donc
\[
  \int_{0}^{\pi} x\sin x \dx = \pi .
\]
Le choix inverse — $u = \sin x$ et $v' = x$ — aurait conduit à une intégrale
plus compliquée que celle de départ : c'est le signe qu'on s'est trompé de
répartition.
\end{exemple}

\subsection{Le changement de variable affine}

\begin{propriete}[changement de variable affine]
Soit $f$ continue et $F$ une primitive de $f$. Pour $\alpha \neq 0$,
\[
  \int_{a}^{b} f(\alpha x + \beta)\dx
  = \frac{1}{\alpha}\Bigl[F(\alpha x + \beta)\Bigr]_{a}^{b}
  = \frac{1}{\alpha}\int_{\alpha a + \beta}^{\alpha b + \beta} f(t)\dt .
\]
\end{propriete}

\begin{exemple}
$\displaystyle\int_{0}^{1} (2x+1)^{3}\dx
= \frac{1}{2}\left[\frac{(2x+1)^{4}}{4}\right]_{0}^{1}
= \frac{1}{8}\left(81-1\right) = 10 .$
\end{exemple}

\begin{entrainement}[calcul d'intégrales]
\exo[\niveau{1}] Calculer $\displaystyle\int_{1}^{2}\left(3x^{2}-4x+1\right)\dx$.

\exo[\niveau{1}] Calculer $\displaystyle\int_{0}^{2} (3x-1)^{4}\dx$.

\exo[\niveau{2}] Calculer $\displaystyle\int_{0}^{\frac{\pi}{2}} x\cos x \dx$
par une intégration par parties.

\exo[\niveau{3}] Calculer $\displaystyle\int_{0}^{1} \frac{x}{\sqrt{x^{2}+1}}\dx$,
puis $\displaystyle\int_{0}^{\pi} x^{2}\sin x\dx$ par deux intégrations par
parties successives.
\end{entrainement}

\begin{corrige}
\textbf{1.} $\bigl[x^{3}-2x^{2}+x\bigr]_{1}^{2} = 2-0 = 2$.

\textbf{2.} $\dfrac{1}{3}\left[\dfrac{(3x-1)^{5}}{5}\right]_{0}^{2}
= \dfrac{1}{15}\left(5^{5}-(-1)^{5}\right) = \dfrac{3126}{15} = \dfrac{1042}{5}$.

\textbf{3.} $u=x$, $v'=\cos x$ : $\bigl[x\sin x\bigr]_{0}^{\pi/2}
- \int_{0}^{\pi/2}\sin x\dx = \dfrac{\pi}{2} - 1$.

\textbf{4.} La première vaut $\bigl[\sqrt{x^{2}+1}\bigr]_{0}^{1} = \sqrt2 - 1$.
Pour la seconde, $u = x^{2}$, $v' = \sin x$ donne
$\bigl[-x^{2}\cos x\bigr]_{0}^{\pi} + 2\int_{0}^{\pi} x\cos x\dx
= \pi^{2} + 2 \times (-2) = \pi^{2}-4$.
\end{corrige}

% ===========================================================================
\section{Aires, valeur moyenne, volumes}

\begin{theoreme}[interprétation d'une intégrale comme une aire]
Soit $f$ continue et \emph{positive} sur $\intervalle{a}{b}$. Alors
$\displaystyle\int_{a}^{b} f(x)\dx$ est l'aire, exprimée en unités d'aire, du
domaine plan défini par
\[
  a \leq x \leq b \qquad \text{et} \qquad 0 \leq y \leq f(x).
\]
\end{theoreme}

\begin{visualisation}[l'aire sous la courbe]
\begin{center}
\begin{tikzpicture}
\begin{axis}[
  width = 10.6cm, height = 5.4cm,
  axis lines = middle, xlabel = {$x$},
  xmin = -0.5, xmax = 5.4, ymin = -0.4, ymax = 3.4,
  xtick = {1,4}, xticklabels = {$a$,$b$}, ytick = \empty,
  axis line style = {bmtGris},
  tick label style = {font=\footnotesize, color=bmtGris}, clip = true,
]
  \addplot[bmtIndigoPale, draw=none, fill=bmtIndigo, fill opacity=0.16,
           domain=1:4, samples=100] {0.35*x^2 - 1.6*x + 3.1} \closedcycle;
  \addplot[bmtIndigo, line width=1.3pt, domain=0.2:5.2, samples=140]
    {0.35*x^2 - 1.6*x + 3.1};
  \draw[bmtGris, dotted] (axis cs:1,0) -- (axis cs:1,1.85);
  \draw[bmtGris, dotted] (axis cs:4,0) -- (axis cs:4,2.3);
  \node[bmtIndigo, font=\footnotesize] at (axis cs:2.5,0.72)
    {$\displaystyle\int_{a}^{b} f(x)\,dx$};
\end{axis}
\end{tikzpicture}
\end{center}

La partie teintée est bornée en bas par l'axe des abscisses, en haut par la
courbe, et sur les côtés par les deux verticales. Son aire est exactement le
nombre que produit le calcul $F(b)-F(a)$. Retenez la condition : la fonction
doit rester positive sur l'intervalle. Si elle change de signe, découpez avec
la relation de Chasles et prenez la valeur absolue sur chaque morceau — sinon
les aires se compensent au lieu de s'ajouter.
\end{visualisation}

\begin{remarque}[unité d'aire]
Si le repère a pour unités $1$ cm sur les deux axes, une unité d'aire vaut
$1$ cm$^{2}$. Si l'unité graphique est $2$ cm, une unité d'aire vaut
$2 \times 2 = 4$ cm$^{2}$ : n'oubliez pas ce facteur, il est régulièrement
demandé dans les sujets.
\end{remarque}

\begin{definition}[valeur moyenne]
La \textbf{valeur moyenne} de $f$ sur $\intervalle{a}{b}$, avec $a<b$, est le
réel
\[
  \mu = \frac{1}{b-a}\int_{a}^{b} f(x)\dx .
\]
\end{definition}

\begin{remarque}
Géométriquement, $\mu$ est la hauteur du rectangle de base $\intervalle{a}{b}$
qui aurait la même aire que le domaine sous la courbe. C'est la traduction
exacte de l'idée de « hauteur moyenne ».
\end{remarque}

\begin{theoreme}[volume d'un solide de révolution]
Soit $f$ continue et positive sur $\intervalle{a}{b}$. Le solide engendré par
la rotation, autour de l'axe des abscisses, du domaine situé sous la courbe de
$f$ a pour volume
\[
  V = \pi \int_{a}^{b} \bigl[f(x)\bigr]^{2}\dx ,
\]
exprimé en unités de volume.
\end{theoreme}

\begin{exemple}[retrouver le volume d'un cône]
Prenons $f(x) = \dfrac{R}{h}x$ sur $\intervalle{0}{h}$, avec $R>0$ et $h>0$.
La rotation de ce segment autour de l'axe des abscisses engendre un cône de
rayon $R$ et de hauteur $h$. Son volume vaut
\[
  V = \pi \int_{0}^{h} \frac{R^{2}}{h^{2}}x^{2}\dx
    = \pi \frac{R^{2}}{h^{2}}\left[\frac{x^{3}}{3}\right]_{0}^{h}
    = \pi \frac{R^{2}}{h^{2}} \times \frac{h^{3}}{3}
    = \frac{\pi R^{2} h}{3},
\]
c'est-à-dire la formule apprise au collège. Le calcul intégral la démontre au
lieu de l'admettre.
\end{exemple}

\begin{entrainement}[aires et volumes]
\exo[\niveau{1}] Calculer l'aire du domaine délimité par la courbe de
$f(x) = x^{2}$, l'axe des abscisses et les droites $x=0$ et $x=3$.

\exo[\niveau{2}] Calculer la valeur moyenne de $x \mapsto x^{2}$ sur
$\intervalle{0}{3}$, et vérifier qu'elle est comprise entre le minimum et le
maximum de la fonction sur cet intervalle.

\exo[\niveau{2}] Calculer l'aire comprise entre les courbes de
$f(x) = x^{2}$ et $g(x) = x$ sur $\intervalle{0}{1}$.

\exo[\niveau{3}] Le domaine sous la courbe de $f(x) = \sqrt{x}$ sur
$\intervalle{0}{4}$ tourne autour de l'axe des abscisses. Calculer le volume
engendré.
\end{entrainement}

\begin{corrige}
\textbf{1.} $\left[\dfrac{x^{3}}{3}\right]_{0}^{3} = 9$ unités d'aire.

\textbf{2.} $\mu = \dfrac{1}{3}\times 9 = 3$, comprise entre $0$ et $9$.

\textbf{3.} Sur $\intervalle{0}{1}$, $x^{2} \leq x$, donc l'aire vaut
$\displaystyle\int_{0}^{1}(x-x^{2})\dx = \dfrac12 - \dfrac13 = \dfrac16$.

\textbf{4.} $V = \pi\displaystyle\int_{0}^{4} x\dx = \pi\left[\dfrac{x^{2}}{2}\right]_{0}^{4} = 8\pi$.
\end{corrige}

% ===========================================================================
% ===========================================================================
\begin{annales}
Les énoncés rassemblés ici sont repris de sujets réellement posés ; leur
provenance est indiquée à chaque fois. Aucun ne fait appel au logarithme ni à
l'exponentielle : tout s'y calcule avec les primitives établies dans ce
chapitre.

\exoann{Baccalauréat, Côte d'Ivoire, séries C et D, session 2023 — exercice 2}
Soit $f$ la fonction définie sur $\R$ par $f(x) = \cos x - x\sin x$. Déterminer
une primitive de $f$ sur $\R$.

\exoann{Baccalauréat, Côte d'Ivoire, séries C et D, session 2023 — exercice 2, suite}
En déduire la valeur de
$\displaystyle\int_{0}^{\frac{\pi}{2}}\left(\cos x - x\sin x\right)\dx$.

\exoann{Concours d'entrée, ENI Fianarantsoa, session 2008-2009 — exercice 3}
Vérifier que pour tout réel $x \neq -1$,
\[
  \frac{x+3}{(x+1)^{3}} = \frac{1}{(x+1)^{2}} + \frac{2}{(x+1)^{3}} .
\]

\exoann{Concours d'entrée, ENI Fianarantsoa, session 2008-2009 — exercice 3, suite}
En déduire la valeur de
$\displaystyle J = \int_{1}^{3}\frac{x+3}{(x+1)^{3}}\dx$.

\medskip
\bmtsource{Les autres énoncés d'annales portant sur le calcul intégral font
intervenir le logarithme ou l'exponentielle. Ils figurent donc aux
chapitres 5 et 6, où ils sont signalés par la mention « calcul intégral ».}
\end{annales}

\begin{corrige}[annales]
\textbf{1.} Le produit $x\cos x$ se dérive en $\cos x - x\sin x$, qui est
exactement $f(x)$. Les primitives de $f$ sur $\R$ sont donc les fonctions
$F(x) = x\cos x + k$, $k\in\R$. Rien n'oblige ici à une intégration par
parties : reconnaître la dérivée d'un produit suffit, et c'est le réflexe que
le sujet cherche à faire jouer.

\textbf{2.} Avec $F(x) = x\cos x$,
\[
  \int_{0}^{\pi/2}\left(\cos x - x\sin x\right)\dx
  = \left[x\cos x\right]_{0}^{\pi/2}
  = \frac{\pi}{2}\cos\frac{\pi}{2} - 0\times\cos 0
  = 0 .
\]
La valeur nulle n'a rien de surprenant : sur $\intervalle{0}{\frac{\pi}{2}}$
la fonction $f$ est positive au début, négative ensuite, et les deux aires se
compensent exactement.

\textbf{3.} On écrit le numérateur en faisant apparaître le dénominateur :
$x+3 = (x+1)+2$. En divisant par $(x+1)^{3}$,
\[
  \frac{x+3}{(x+1)^{3}} = \frac{x+1}{(x+1)^{3}} + \frac{2}{(x+1)^{3}}
  = \frac{1}{(x+1)^{2}} + \frac{2}{(x+1)^{3}} .
\]

\textbf{4.} Chacun des deux morceaux est de la forme $u'u^{n}$ avec
$u = x+1$ :
\[
  \int_{1}^{3}\frac{\dx}{(x+1)^{2}} = \left[\frac{-1}{x+1}\right]_{1}^{3}
  = -\frac14+\frac12 = \frac14,
  \qquad
  \int_{1}^{3}\frac{2\dx}{(x+1)^{3}} = \left[\frac{-1}{(x+1)^{2}}\right]_{1}^{3}
  = -\frac{1}{16}+\frac14 = \frac{3}{16}.
\]
En additionnant, $J = \dfrac14+\dfrac{3}{16} = \dfrac{7}{16}$.
\end{corrige}

\begin{jecherche}[le réservoir de la briqueterie]
Une briqueterie d'Antsirabe possède un réservoir obtenu en faisant tourner
autour de l'axe des abscisses la courbe de
\[
  f(x) = \sqrt{4-\frac{x^{2}}{4}}, \qquad x \in \intervalle{0}{4},
\]
les longueurs étant exprimées en mètres.

\begin{enumerate}[label=\textbf{\arabic*.}, leftmargin=*, itemsep=0.3em]
  \item Vérifier que $f$ est bien définie et positive sur $\intervalle{0}{4}$,
        et calculer $f(0)$ et $f(4)$.
  \item Quelle forme géométrique la courbe de $f$ décrit-elle ? En déduire
        l'allure du réservoir.
  \item Calculer le volume du réservoir en mètres cubes.
  \item Le réservoir est rempli jusqu'à l'abscisse $x=2$. Quel volume d'eau
        contient-il alors ? Quelle fraction de sa capacité cela représente-t-il ?
  \item L'usine veut repeindre la surface du disque de section à l'abscisse
        $x$. Exprimer cette aire en fonction de $x$, puis déterminer sa valeur
        moyenne sur $\intervalle{0}{4}$.
\end{enumerate}
\end{jecherche}

\begin{corrige}[le réservoir de la briqueterie]
\textbf{1.} $4-\dfrac{x^{2}}{4} \geq 0$ équivaut à $x^{2} \leq 16$, vrai sur
$\intervalle{0}{4}$ : $f$ y est définie et positive. De plus $f(0) = 2$ et
$f(4) = 0$.

\textbf{2.} En posant $y = f(x)$ et en élevant au carré,
$y^{2} = 4-\dfrac{x^{2}}{4}$, c'est-à-dire
$\dfrac{x^{2}}{16}+\dfrac{y^{2}}{4} = 1$ : la courbe est un quart d'ellipse. Le
réservoir est donc un demi-ellipsoïde de révolution, large de $2$ m au fond et
effilé vers la sortie.

\textbf{3.} $V = \pi\displaystyle\int_{0}^{4}\left(4-\frac{x^{2}}{4}\right)\dx
= \pi\left[4x-\frac{x^{3}}{12}\right]_{0}^{4}
= \pi\left(16-\frac{16}{3}\right) = \frac{32\pi}{3} \approx 33{,}5$ m$^{3}$.

\textbf{4.} $\pi\displaystyle\int_{0}^{2}\left(4-\frac{x^{2}}{4}\right)\dx
= \pi\left(8-\frac{2}{3}\right) = \frac{22\pi}{3} \approx 23{,}0$ m$^{3}$, soit
$\dfrac{22}{32} = \dfrac{11}{16}$, environ $69\,\%$ de la capacité. Remplir la
moitié de la longueur remplit donc bien plus que la moitié du volume : la
partie large est au fond.

\textbf{5.} L'aire du disque de section vaut
$\pi\left[f(x)\right]^{2} = \pi\left(4-\dfrac{x^{2}}{4}\right)$. Sa valeur
moyenne sur $\intervalle{0}{4}$ est
$\dfrac{1}{4}\times\dfrac{32\pi}{3} = \dfrac{8\pi}{3} \approx 8{,}4$ m$^{2}$.
\end{corrige}


% ===========================================================================
\begin{jemeteste}
Répondre par vrai ou faux, en justifiant en une ligne.

\begin{enumerate}[label=\textbf{\arabic*.}, leftmargin=*, itemsep=0.28em]
  \item Une fonction continue sur un intervalle y admet une unique primitive.
  \item $\displaystyle\int_{a}^{b} f(x)\dx$ est toujours positive.
  \item Si $f \leq g$ sur $\intervalle{a}{b}$ avec $a<b$, alors
        $\int_{a}^{b} f \leq \int_{a}^{b} g$.
  \item L'intégration par parties transforme toujours l'intégrale en une plus
        simple.
  \item La valeur moyenne d'une fonction sur un intervalle est comprise entre
        son minimum et son maximum sur cet intervalle.
  \item Le volume d'un solide de révolution vaut $\int_{a}^{b} \pi f(x)\dx$.
\end{enumerate}
\end{jemeteste}

\begin{corrige}[je me teste]
\textbf{1.} Faux — elle en admet une infinité, qui diffèrent d'une constante.
\textbf{2.} Faux — cela dépend du signe de $f$ et de l'ordre des bornes.
\textbf{3.} Vrai.
\textbf{4.} Faux — seulement si $u$ et $v'$ sont bien choisis ; le mauvais
choix complique l'intégrale.
\textbf{5.} Vrai — conséquence de l'encadrement $m(b-a) \leq \int f \leq M(b-a)$.
\textbf{6.} Faux — c'est $\pi\int_{a}^{b} \left[f(x)\right]^{2}\dx$ : le carré
est indispensable.
\end{corrige}

% ===========================================================================
\begin{commeaubac}
\bmtsource{Exercice au format de l'épreuve}
Soit $f$ la fonction définie sur $\intervalle{0}{2}$ par
$f(x) = x\sqrt{4-x^{2}}$.

\begin{enumerate}[label=\textbf{\arabic*.}, leftmargin=*, itemsep=0.25em]
  \item Vérifier que $f$ est positive sur $\intervalle{0}{2}$ et calculer
        $f(0)$ et $f(2)$. \hfill\textup{(0,5 pt)}
  \item Calculer $f'(x)$ et dresser le tableau de variation de $f$ sur
        $\intervalle{0}{2}$. \hfill\textup{(1,5 pt)}
  \item Déterminer une primitive de $f$ sur $\intervalle{0}{2}$.
        \emph{On reconnaîtra une forme $u'u^{n}$.} \hfill\textup{(1 pt)}
  \item En déduire l'aire, en unités d'aire, du domaine délimité par la courbe
        de $f$, l'axe des abscisses et les droites $x=0$ et $x=2$.
        \hfill\textup{(1 pt)}
  \item Calculer la valeur moyenne de $f$ sur $\intervalle{0}{2}$.
        \hfill\textup{(0,5 pt)}
  \item Calculer le volume du solide engendré par la rotation de ce domaine
        autour de l'axe des abscisses. \hfill\textup{(1,5 pt)}
\end{enumerate}
\end{commeaubac}

\begin{corrige}[exercice au format de l'épreuve]
\textbf{1.} Sur $\intervalle{0}{2}$, $x \geq 0$ et $4-x^{2} \geq 0$ : $f \geq 0$.
$f(0) = f(2) = 0$.

\textbf{2.} $f'(x) = \sqrt{4-x^{2}} - \dfrac{x^{2}}{\sqrt{4-x^{2}}}
= \dfrac{4-2x^{2}}{\sqrt{4-x^{2}}}$, du signe de $2-x^{2}$ : $f$ croît sur
$\intervalle{0}{\sqrt2}$ puis décroît sur $\intervalle{\sqrt2}{2}$, avec un
maximum $f(\sqrt2) = 2$.

\textbf{3.} $f(x) = -\dfrac{1}{2}\times(-2x)\left(4-x^{2}\right)^{1/2}$, forme
$-\frac12 u'u^{1/2}$ avec $u = 4-x^{2}$ : une primitive est
$-\dfrac{1}{3}\left(4-x^{2}\right)^{3/2}$.

\textbf{4.} $\left[-\dfrac{(4-x^{2})^{3/2}}{3}\right]_{0}^{2} = 0 + \dfrac{8}{3}
= \dfrac{8}{3}$ unités d'aire.

\textbf{5.} $\mu = \dfrac{1}{2}\times\dfrac{8}{3} = \dfrac{4}{3}$.

\textbf{6.} $V = \pi\displaystyle\int_{0}^{2} x^{2}\left(4-x^{2}\right)\dx
= \pi\left[\dfrac{4x^{3}}{3}-\dfrac{x^{5}}{5}\right]_{0}^{2}
= \pi\left(\dfrac{32}{3}-\dfrac{32}{5}\right) = \dfrac{64\pi}{15}$.
\end{corrige}

% ===========================================================================
\begin{concours}
\bmtsource{Concours d'entrée, ENI Fianarantsoa, session 2002-2003 — exercice 4}
Pour tout entier naturel $n$, on pose
$\displaystyle I_{n} = \int_{0}^{\frac{\pi}{2}} \sin^{n}x \dx$.

\begin{enumerate}[label=\textbf{\arabic*.}, leftmargin=*, itemsep=0.25em]
  \item Calculer $I_{0}$ et $I_{1}$.
  \item Montrer que pour tout entier $n \geq 2$,
        $n\,I_{n} = (n-1)\,I_{n-2}$.
  \item En déduire les valeurs de $I_{2}$, $I_{4}$ et $I_{6}$.
\end{enumerate}
\end{concours}

\begin{corrige}[vers le concours]
\textbf{1.} $I_{0} = \displaystyle\int_{0}^{\pi/2}\dx = \frac{\pi}{2}$ et
$I_{1} = \left[-\cos x\right]_{0}^{\pi/2} = 1$.

\textbf{2.} Pour $n \geq 2$, on écrit $\sin^{n}x = \sin^{n-1}x \times \sin x$
et on intègre par parties avec $u = \sin^{n-1}x$ et $v' = \sin x$, donc
$u' = (n-1)\sin^{n-2}x\cos x$ et $v = -\cos x$ :
\[
  I_{n} = \left[-\cos x \sin^{n-1}x\right]_{0}^{\pi/2}
        + (n-1)\int_{0}^{\pi/2}\cos^{2}\!x\,\sin^{n-2}x \dx .
\]
Le crochet est nul aux deux bornes. En remplaçant $\cos^{2}x$ par
$1-\sin^{2}x$, l'intégrale restante se sépare en $I_{n-2}-I_{n}$, d'où
\[
  I_{n} = (n-1)I_{n-2} - (n-1)I_{n},
  \qquad\text{c'est-à-dire}\qquad
  n\,I_{n} = (n-1)\,I_{n-2}.
\]

\textbf{3.} La relation $I_{n} = \dfrac{n-1}{n}I_{n-2}$ se déroule à partir de
$I_{0}$ :
\[
  I_{2} = \frac12 \times \frac{\pi}{2} = \frac{\pi}{4},
  \qquad
  I_{4} = \frac34 \times \frac{\pi}{4} = \frac{3\pi}{16},
  \qquad
  I_{6} = \frac56 \times \frac{3\pi}{16} = \frac{5\pi}{32}.
\]
\end{corrige}


% ===========================================================================
\begin{devoir}[2 heures]
Trois exercices indépendants, notés sur $20$. Le devoir porte sur ce chapitre
et sur les trois qui le précèdent. Les énoncés sont repris de sessions réelles
et assemblés ici en épreuve ; leur provenance est indiquée à chaque fois.

\exods{1}{D'après le baccalauréat probatoire, Cameroun, série C, session 2025 — exercice 2}{5 points}
Soit $f$ la fonction définie par $f(x) = \dfrac{x^{2}-2x+2}{x-1}$.
\begin{enumerate}[label=\textbf{\arabic*.}, leftmargin=*, itemsep=0.2em]
  \item Montrer que $f(x) = x-1+\dfrac{1}{x-1}$ pour tout $x \neq 1$.
        \hfill\textup{(1 pt)}
  \item Calculer les limites de $f$ aux bornes de son ensemble de définition
        et en déduire les deux asymptotes à sa courbe. \hfill\textup{(2 pts)}
  \item Vérifier que $f'(x) = \dfrac{x^{2}-2x}{(x-1)^{2}}$, puis dresser le
        tableau de variation de $f$ sur $\intervalleo{1}{+\infty}$.
        \hfill\textup{(2 pts)}
\end{enumerate}

\exods{2}{D'après le concours ENI Fianarantsoa, session 2008-2009 (exercice 3) et le baccalauréat, Côte d'Ivoire, séries C et D, session 2023 (exercice 2)}{7 points}
\textbf{Partie A.}
\begin{enumerate}[label=\textbf{\arabic*.}, leftmargin=*, itemsep=0.2em]
  \item Vérifier que pour tout réel $x \neq -1$,
        $\dfrac{x+3}{(x+1)^{3}} = \dfrac{1}{(x+1)^{2}}+\dfrac{2}{(x+1)^{3}}$.
        \hfill\textup{(1,5 pt)}
  \item En déduire la valeur de
        $\displaystyle J = \int_{1}^{3}\frac{x+3}{(x+1)^{3}}\dx$.
        \hfill\textup{(2 pts)}
\end{enumerate}
\textbf{Partie B.} Soit $g$ la fonction définie sur $\R$ par
$g(x) = \cos x-x\sin x$.
\begin{enumerate}[label=\textbf{\arabic*.}, leftmargin=*, itemsep=0.2em, resume]
  \item Déterminer une primitive de $g$ sur $\R$. \hfill\textup{(1,5 pt)}
  \item Calculer
        $\displaystyle\int_{0}^{\frac{\pi}{2}}\left(\cos x-x\sin x\right)\dx$.
        \hfill\textup{(1 pt)}
  \item Le résultat est nul, alors que $g$ n'est pas la fonction nulle.
        Expliquer ce que cela signifie pour la courbe de $g$ sur
        $\intervalle{0}{\frac{\pi}{2}}$. \hfill\textup{(1 pt)}
\end{enumerate}

\exods{3}{D'après le concours d'entrée, ENI Fianarantsoa, session 2002-2003 — exercice 4}{8 points}
Pour tout entier naturel $n$, on pose
$\displaystyle I_{n} = \int_{0}^{\frac{\pi}{2}}\sin^{n}x\dx$.
\begin{enumerate}[label=\textbf{\arabic*.}, leftmargin=*, itemsep=0.2em]
  \item Calculer $I_{0}$ et $I_{1}$. \hfill\textup{(1,5 pt)}
  \item Démontrer que pour tout entier $n \geq 2$,
        $n\,I_{n} = (n-1)\,I_{n-2}$. \hfill\textup{(3 pts)}
  \item En déduire les valeurs exactes de $I_{2}$, $I_{4}$ et $I_{6}$.
        \hfill\textup{(1,5 pt)}
  \item Montrer que la suite $(I_{n})$ est positive et décroissante, puis
        conclure quant à sa convergence. \hfill\textup{(2 pts)}
\end{enumerate}
\end{devoir}

\begin{corrige}[devoir surveillé]
\textbf{Exercice 1. 1.} $x^{2}-2x+2 = (x-1)^{2}+1$, donc
$f(x) = \dfrac{(x-1)^{2}+1}{x-1} = x-1+\dfrac{1}{x-1}$.

\textbf{2.} $\Df = \R\setminus\{1\}$. En $1^{-}$, $\frac{1}{x-1}\to-\infty$
donc $f(x)\to-\infty$ ; en $1^{+}$, $f(x)\to+\infty$ : la droite $x = 1$ est
asymptote verticale. Aux infinis, $f(x)-(x-1) = \frac{1}{x-1}\to0$ : la droite
$y = x-1$ est asymptote oblique.

\textbf{3.} $f'(x) = 1-\dfrac{1}{(x-1)^{2}} = \dfrac{(x-1)^{2}-1}{(x-1)^{2}}
= \dfrac{x^{2}-2x}{(x-1)^{2}}$, du signe de $x(x-2)$. Sur
$\intervalleo{1}{+\infty}$, ce produit est négatif avant $2$ et positif
après : $f$ décroît de $+\infty$ jusqu'à $f(2) = 2$, puis croît vers
$+\infty$.

\medskip
\textbf{Exercice 2. 1.} On écrit le numérateur en faisant apparaître le
dénominateur : $x+3 = (x+1)+2$, d'où
$\dfrac{x+3}{(x+1)^{3}} = \dfrac{x+1}{(x+1)^{3}}+\dfrac{2}{(x+1)^{3}}
= \dfrac{1}{(x+1)^{2}}+\dfrac{2}{(x+1)^{3}}$.

\textbf{2.} Chaque morceau est de la forme $u'u^{n}$ avec $u = x+1$ :
\[
  \int_{1}^{3}\frac{\dx}{(x+1)^{2}} = \left[\frac{-1}{x+1}\right]_{1}^{3}
  = \frac14, \qquad
  \int_{1}^{3}\frac{2\dx}{(x+1)^{3}} = \left[\frac{-1}{(x+1)^{2}}\right]_{1}^{3}
  = \frac{3}{16},
\]
d'où $J = \dfrac14+\dfrac{3}{16} = \dfrac{7}{16}$.

\textbf{3.} La dérivée du produit $x\cos x$ vaut $\cos x-x\sin x$, c'est-à-dire
$g(x)$ : les primitives de $g$ sont les fonctions $G(x) = x\cos x+k$. Une
intégration par parties n'est pas nécessaire — reconnaître la dérivée d'un
produit suffit.

\textbf{4.} $\displaystyle\int_{0}^{\pi/2}g(x)\dx
= \left[x\cos x\right]_{0}^{\pi/2}
= \frac{\pi}{2}\cos\frac{\pi}{2}-0 = 0$.

\textbf{5.} La fonction $g$ n'est pas nulle : elle vaut $1$ en $0$ et
$-\frac{\pi}{2}$ en $\frac{\pi}{2}$. Elle change donc de signe sur
l'intervalle, et l'intégrale nulle signifie que l'aire comprise entre la
courbe et l'axe des abscisses au-dessus de l'axe est exactement égale à celle
située au-dessous : les deux se compensent.

\medskip
\textbf{Exercice 3. 1.} $I_{0} = \displaystyle\int_{0}^{\pi/2}\dx
= \frac{\pi}{2}$ et
$I_{1} = \left[-\cos x\right]_{0}^{\pi/2} = 0+1 = 1$.

\textbf{2.} Pour $n \geq 2$, on écrit $\sin^{n}x = \sin^{n-1}x \times \sin x$
et l'on intègre par parties avec $u = \sin^{n-1}x$ et $v' = \sin x$, donc
$u' = (n-1)\sin^{n-2}x\cos x$ et $v = -\cos x$ :
\[
  I_{n} = \left[-\cos x\,\sin^{n-1}x\right]_{0}^{\pi/2}
  +(n-1)\int_{0}^{\pi/2}\cos^{2}\!x\,\sin^{n-2}x\dx .
\]
Le crochet est nul aux deux bornes. En remplaçant $\cos^{2}x$ par
$1-\sin^{2}x$, l'intégrale restante se sépare en $I_{n-2}-I_{n}$, d'où
$I_{n} = (n-1)I_{n-2}-(n-1)I_{n}$, c'est-à-dire $n\,I_{n} = (n-1)\,I_{n-2}$.

\textbf{3.} La relation $I_{n} = \frac{n-1}{n}I_{n-2}$ se déroule à partir de
$I_{0}$ :
\[
  I_{2} = \frac12\times\frac{\pi}{2} = \frac{\pi}{4}, \qquad
  I_{4} = \frac34\times\frac{\pi}{4} = \frac{3\pi}{16}, \qquad
  I_{6} = \frac56\times\frac{3\pi}{16} = \frac{5\pi}{32}.
\]

\textbf{4.} Sur $\intervalle{0}{\frac{\pi}{2}}$ on a $0 \leq \sin x \leq 1$,
donc $\sin^{n}x \geq 0$ : chaque $I_{n}$ est positif. De plus
$\sin^{n+1}x \leq \sin^{n}x$ sur cet intervalle, et l'intégration conserve
l'inégalité : $I_{n+1} \leq I_{n}$. La suite est donc décroissante et minorée
par $0$ : elle converge. On peut ajouter que sa limite est nulle, puisque
$n\,I_{n} = (n-1)I_{n-2}$ force $I_{n}$ à tendre vers $0$.
\end{corrige}
